
Reward models serve as optimization targets in reinforcement learning from human feedback (RLHF), guiding language model training toward behaviors that human raters prefer. These models are typically assumed to evaluate response content---accuracy, helpfulness, and reasoning quality. This study examines whether RLHF-trained reward models also encode preferences for structural formatting features independent of content.

Prior work on reward model evaluation has focused primarily on content dimensions. RewardBench evaluates reward models on helpfulness, safety, and reasoning accuracy. MRMBench probes multi-dimensional preferences including instruction following. These benchmarks assess whether reward models correctly identify higher-quality responses but do not examine systematic preferences for particular formatting patterns when content quality is controlled.

The present investigation addresses this gap by testing whether reward models assign different scores to responses that present identical content in different structural formats---specifically, enumerated lists versus synthesized prose.

\subsubsection{1.1 Research Questions}

This study addresses three research questions:

\begin{enumerate}
\item Do RLHF-trained reward models exhibit significant enumeration preference under controlled conditions where content, length, and quality are matched?
\end{enumerate}

\begin{enumerate}
\item Does enumeration preference replicate across architecturally distinct reward models?
\end{enumerate}

\begin{enumerate}
\item Does the effect depend on model architecture (decoder-only vs. encoder-only)?
\end{enumerate}

\subsubsection{1.2 Overview of Approach}

To isolate structural effects from content quality, this study employs controlled stimulus pairs where enumerated and synthesized responses contain semantically equivalent information matched within $\pm 2$\% token length. A $2\times 2$ factorial design (correctness $\times$ completeness) enables analysis of potential interactions between structure and quality dimensions. Four reward models spanning different architectures and training objectives are evaluated to assess cross-model generalization.
